% !TEX root = ../main.tex
\section{Conclusions and Recommendations}
\label{sec:conclusions}
\label{sec:recommendations}

% Merged 2026-10-11 with the former Recommendations section per the owner's note
% on the printed report, page 19: "Combine w/ Conclusions."
% (handwritten/2026-10-11_main.md). 05_recommendations.tex is removed; its
% call-out is the second one below, with the phrase the owner struck taken out.

\plantodo{S15, S16}{Briefly restate the problem, approach, and findings, including potential benefits in relative terms if absolute figures are sensitive or incomplete (e.g., ``reduces network-level delay propagation by an illustrative X\% under the demonstrated scenario'' rather than an unsupported absolute claim). Cover: known architecture omissions, limitations of the selected decomposition (S6), limitations of the experiment (Sections~\ref{sec:assumptions} and~\ref{sec:scoring-limits}), and opportunities for future expansion (e.g., additional off-nominal scenarios, unexplored S9 conflicts, deeper crew/maintenance-routing integration).}

\plantodo{S15, S16 (step 10, 2026-11-10)}{\textbf{Close the two claims.} Claim 1: say whether the higher-fidelity cases resolved the conflict better as measured against the kept objectives (Section~\ref{sec:experiment-results}), for one conflict in one sector, and what that does and does not support. Claim 2: the model flagged the gap from the architecture alone, before any simulation (Section~\ref{sec:model-gap}); the requirement breakdown set out the means of closing it; the analysis confirmed and measured it; and the fidelity was swapped with the structure unchanged. The gap was already known from the literature; the claim is that modeling could have found it proactively (D-044). Say what the reading of D-044 gives up: no requirement fails on decision quality, so B minus A is reported, not graded.}

\plantodo{S14, S15}{\textbf{Recommendations.} Restate the main findings from Section~\ref{sec:results} and provide recommendations for moving forward. Include the alternatives considered (alternative stakeholder-weighting schemes from S11) and the criteria used to select the recommended approach. Suggest additional areas of study (candidate: architecture extensions, additional off-nominal scenarios from S15, unexplored myopic-optimization conflicts from S9's catalog).}
